\documentclass[11pt]{article}
\usepackage[T1]{fontenc}
\usepackage{lmodern}
\usepackage[margin=1in,top=.8in,bottom=.85in]{geometry}
\usepackage{amsmath,amssymb,amsthm,booktabs,microtype}
\usepackage[hidelinks]{hyperref}
\newtheorem{theorem}{Theorem}
\newtheorem{lemma}[theorem]{Lemma}
\newtheorem{corollary}[theorem]{Corollary}
\newcommand{\N}{\mathbb N}
\newcommand{\Z}{\mathbb Z}
\newcommand{\norm}[1]{\left|#1\right|}
\title{Digit-sum residues on arithmetic progressions\\in every rational base}
\author{Blake Stone mathematics repository\\[-2pt]\small Proof developed with OpenAI Codex}
\date{November 26, 2026}
\begin{document}
\maketitle
\begin{abstract}
Cassaigne, Espinoza, Rigo and Stipulanti asked whether the sum of the digits in
base $3/2$ is equally distributed modulo every $m$. We prove this for every
arithmetic progression and classify the corresponding distribution for every reduced
rational base $p/q>1$. For starts $Aj+b$, put $g=\gcd(m,A,p-q)$. The limiting
frequency of a residue $c$ is $g/m$ if $c\equiv qb\pmod g$ and zero otherwise.
Convergence is uniform in the starting point for each fixed step. The proof pairs
progression terms whose higher digits agree and whose lower digit sums give
different roots of unity.
\end{abstract}

\section{The distribution}
Fix coprime integers $p>q\ge1$. We use the rational-base digits of Akiyama,
Frougny and Sakarovitch \cite{afs}, with the division algorithm and normalization
used in \cite{mst}. Set
\[
 T(n)=\lfloor qn/p\rfloor,\qquad d(n)=qn-pT(n),\qquad
 S(0)=0,\qquad S(n)=d(n)+S(T(n)).
\]
Here $0\le d(n)<p$ and $T(n)<n$ for positive $n$. At $p/q=3/2$, reducing
$S$ modulo $m$ sends each parent pair $ab$ to $a,(a+2),(b+1)$ modulo $m$.
This is precisely the word in the concluding question of \cite{cers}.
Their Theorem~18 treats the two individual symbols modulo two along dyadic
residue classes. Morgenbesser, Steiner and Thuswaldner studied digit patterns
within the representations and the summatory digit sum \cite{mst}.

\begin{theorem}\label{main}
For $m\ge2$, $A\ge1$, $b\ge0$, and $0\le c<m$, let $g=\gcd(m,A,p-q)$.
Then
\[
\lim_{N\to\infty}\frac1N\#\{0\le j<N:S(Aj+b)\equiv c\pmod m\}
=\begin{cases} g/m,&c\equiv qb\pmod g,\\ 0,&\text{otherwise}.\end{cases}
\]
For fixed $p,q,m,A$ the error tends to zero uniformly over $b\ge0$.
\end{theorem}
For example, base $3/2$ gives frequency $1/m$ for every residue on every
progression. In base $5/2$, the sums at $6j+1$ modulo $6$ occupy just
residues $2$ and $5$, each with frequency $1/2$. The restriction follows
already from
\begin{equation}
 S(n)\equiv qn\pmod{p-q}.\label{congruence}
\end{equation}
Indeed $d(n)\equiv q(n-T(n))\pmod{p-q}$, and successive terms telescope.

\section{One carry and a positive proportion of cancellations}
Write $F_k=T^k$, $P=p^k$, $Q=q^k$ and
$L_k(n)=\sum_{i=0}^{k-1}d(T^i(n))$. Then
\begin{align}
 S(n)&=L_k(n)+S(F_k(n)),\nonumber\\
 F_k(Pj+n)&=Qj+F_k(n),\qquad L_k(Pj+n)=L_k(n).\label{shift}
\end{align}
The identities follow by induction from $T(pj+n)=qj+T(n)$ and $d(pj+n)=d(n)$.
The map $F_k$ is nondecreasing. Put $e=p-q$ and $\Delta(n)=S(n+1)-S(n)$.
When $d(n)+q<p$ the parent stays fixed and $\Delta(n)=q$. Otherwise it
advances one place and $\Delta(n)=q-p+\Delta(T(n))$.

\begin{lemma}\label{witness}
If $\gcd(A,p)=1$, there is a class $c_0$ modulo $p^2$ such that
$\Delta(n)=q$ and $\Delta(n+A)=q-e$ for every $n\equiv c_0\pmod{p^2}$.
\end{lemma}
\begin{proof}
Let $v$ be the least residue of $qA$ modulo $p$, so $1\le v<p$.
Set $l=\max(0,p-q-v)$ and $z=l+v$. Then $l<p-q\le z<p$.
Choose $c_0$ with $q(c_0+A)\equiv z\pmod{p^2}$ using $\gcd(q,p^2)=1$.
At $n$ in this class, $d(n)=l$ and $d(n+A)=z$. The parent $T(n+A)$
is divisible by $p$, so its digit is zero. These give respectively zero
carries at $n$ and exactly one at $n+A$, proving the identities.
\end{proof}

Fix an integer $h\ge2$ with $h\nmid e$. Call a residue $n$ modulo $P$
\emph{marked} if $F_k(n+A)=F_k(n)$ and
$L_k(n+A)\not\equiv L_k(n)\pmod h$. Both conditions are periodic by
\eqref{shift}. Write $I(n)$ for the marking indicator and
$D(n)=F_k(n+A)-F_k(n)\ge0$ for the parent displacement.
\begin{lemma}\label{mass}
For every $A$ coprime to $p$, some $k\ge2$ satisfies
$\sum_{r=0}^{P-1}I(Ar+b)\ge P/(4p^2)$ for every $b$.
\end{lemma}
\begin{proof}
The sum of $D$ over a period is $AQ$. To see this, expand each difference
as the sum of $A$ successive increments of $F_k$. For each shift the
increments over a whole period telescope to $Q$. If $D(n)+I(n)=0$,
then $V(n)=S(n+A)-S(n)$ is zero modulo $h$. But at the indices of
Lemma~\ref{witness} one has $V(n+1)-V(n)=-e$, which is nonzero modulo $h$.
At least one of $D(n),I(n),D(n+1),I(n+1)$ is positive there. Sum this
inequality over a period. For $k\ge2$ it gives
\[
 P/p^2\le 2AQ+2\sum_{n=0}^{P-1} I(n).
\]
Choose $k$ large enough that $4AQ\le P/p^2$, using $q/p<1$.
Multiplication by $A$ followed by addition of $b$ permutes the classes modulo $P$.
\end{proof}

\section{Cancellation for all progression steps}
Let $\zeta$ be a root of unity of order $h\nmid e$ and put $u(n)=\zeta^{S(n)}$.
Since two distinct $h$th roots of unity have a sum of modulus strictly
less than two, there is $\eta\in(0,1]$ such that the modulus of every such
sum is at most $2-\eta$. For $N\ge1$ set
\[
 M_A(N)=\sup_{b\ge0}\left|\frac1N\sum_{j=0}^{N-1}u(Aj+b)\right|,\qquad
 B=\sup_{\gcd(A,p)=1,\ A\ge1}\limsup_{N\to\infty}M_A(N).
\]
The order of the supremum over steps and the upper limit is important.
Every $M_A(N)$ is at most one.

Fix an admissible $A$ and take $k$ from Lemma~\ref{mass}. In a block of
$PH$ indices write $j=Pt+r$ and $n_r=Ar+b$. By \eqref{shift},
\begin{equation}
 u(A(Pt+r)+b)=\zeta^{L_k(n_r)}u(AQt+F_k(n_r)).\label{branch}
\end{equation}
Use the sum of two copies of the block, the second shifted one progression
index to the right. Their difference comprises two endpoints of modulus one.
For each $r$, pair its branch \eqref{branch} with the branch at $n_r+A$.
When $n_r$ is marked the two parent offsets coincide, so the sum over $t<H$
has modulus at most $(2-\eta)H M_{AQ}(H)$. At every other $r$ the triangle
inequality bounds it by $2H M_{AQ}(H)$. This also applies at $r=P-1$,
since \eqref{branch} accepts any offset. Adding the bounds gives
\begin{equation}
 M_A(PH)\le \left(1-\frac{\eta}{8p^2}\right) M_{AQ}(H)+\frac1{PH}.
 \label{contraction}
\end{equation}
For arbitrary $N=PH+R$ with $0\le R<P$, deleting the last $R$ terms changes
the normalized sum by at most $2P/N$ when $H>0$. Since $AQ$ is coprime to
$p$, taking upper limits for each fixed $A$, then the supremum over $A$,
gives $B\le(1-\eta/(8p^2))B$. Thus $B=0$.

For a general $A$, write $A=UV$ where the prime factors of $U$ divide $p$
and $\gcd(V,p)=1$. Choose $s$ with $U\mid p^s$. Splitting $j$ into its
$p^s/U$ residue classes puts each branch of $Aj+b$ in the form
$p^s Vt+n_r$. Equation~\eqref{shift} gives a fixed unit coefficient
and parent step $q^s V$, which is coprime to $p$. Each branch mean tends
to zero uniformly in its offset. The finitely many branches and their
at most one-term length differences give $M_A(N)\to0$.

\section{Recovering the residue frequencies}
Set $\omega=\exp(2\pi i/m)$. The indicator of $S(n)\equiv c\pmod m$
is $m^{-1}\sum_{t=0}^{m-1}\omega^{t(S(n)-c)}$. Terms for which the
order of $\omega^t$ does not divide $e$ have uniformly vanishing
progression means by Section~3. For each remaining character,
\eqref{congruence} replaces $S(Aj+b)$ by $q(Aj+b)$.
This geometric progression has mean zero unless the order divides
$qA$, in which case its value is constantly $\omega^{t(qb-c)}$.
The geometric-sum bound for each nonconstant progression is independent
of $b$. Since $\gcd(e,q)=1$, the surviving characters are precisely
the $g=\gcd(m,A,e)$ characters with order dividing $g$.
Orthogonality on $\Z/g\Z$ now gives Theorem~\ref{main}.

\paragraph{Connection with binary factor frequencies.}
At $p/q=3/2$ and $m=2$, Theorem~\ref{main} supplies the digit-sum parity
limit on every dyadic progression. Feeding this result into the finite-word
recursion of \cite{factor} proves the exact binary factor frequencies and
their complement and reversal symmetries directly from the sequence definitions.
The accompanying Lean development checks Theorem~\ref{main} for arbitrary
parameters and this specialization.

\begin{thebibliography}{9}
\bibitem{afs} S. Akiyama, C. Frougny and J. Sakarovitch, \emph{Powers of rationals
modulo 1 and rational base number systems}, Israel J. Math. \textbf{168} (2008),
53--91.
\bibitem{cers} J. Cassaigne, B. Espinoza, M. Rigo, and M. Stipulanti,
\emph{Symbols frequencies in the Thue--Morse word in base $3/2$ and related conjectures},
arXiv:2602.21895 (2026), Theorem~18 and Section~4.
\url{https://arxiv.org/abs/2602.21895}.
\bibitem{mst} J. F. Morgenbesser, W. Steiner, and J. M. Thuswaldner,
\emph{Patterns in rational base number systems}, J. Fourier Anal. Appl.
\textbf{19} (2013), 225--250. \url{https://doi.org/10.1007/s00041-012-9246-1}.
\bibitem{factor} Blake Stone mathematics repository,
\emph{Factor frequencies and fifth powers in the base-$3/2$ Thue--Morse words}
(2026). \href{https://github.com/blakestone-x/math/tree/9abc19936fb4297de972941c1aad6b4659da5610/Dekking-base-three-halves}{Repository manuscript and Lean development}.
\end{thebibliography}
\end{document}
